% ============================================================
% M3 S3 练习件·W1 臂1 样板件 —— 课时01 坐标法（2.1）＝练习件线母版（后波照抄件）
% 生成：M3 成卷轮 S3 母版代理｜2026-09-14｜编译 cwd＝本目录（中文路径禁 \input 绝对路径）
%   xelatex main-true.tex ×2 → 含详解印本；xelatex main-false.tex ×2 → 纯题版（单源双本）
% 输入链（全只读）：题面库 课时01-坐标法.md（题面逐字装配）＋ -答案侧.md（值逐字回嵌）
%   ＋ 定稿/批A-课时01-坐标法.md（详解回嵌源）；toolchain＝qp-m3.sty（钉后版
%   md5 7c3930362be8a0a2bdf21bbf8ac16573，本地挂载副本，破损即停）。
% 母版体例：详见《练习件母版体例说明.md》；门谱读数见 过程件 工作区/_tmpM3S3母版0914/。
% 键账：件内 ansblock 键集 ≡ 题面库片练键集 ≡ manifest 练键序 E1~E16（16 键，零缺漏零多余）；
%   拓区隔离：2章-拓/测/滚 键不入本件（S4 拓展册收）；导学键 G1~G5 属导学件域，亦不入。
% 卷式例外案B：练习件不属卷式——答案块物理落点恒为题后 ansblock，禁卷末附卷制；
%   全线括线模（导言 \ansblockgrayfalse 一次置定，件内不切换；灰底盒不可跨栏断，练习件禁用）。
% 题序＝manifest 键序 E1~E16（零重排）；印面题号 1~16 连号（\tihao 号＝\ansitem 号同键）；
%   三档切位 1~7／8~14／15~16（照库序切位，非难度重排；15~16＝0913命制入槽新命对）。
% 槽配实录：单选3（E1/E15/E16）＋多选2（E2/E3，≤4 合规）＋填空4（E4~E7）＋解答7（E8~E14）；
%   10选4填2解 为波次方案基线槽配，本片实配照题面侧头标（批A §四 满足度表同口径）。
% 选项槽位：默认四连排（0.25\linewidth−0.25\lxhang），超宽降档梯 四连排→两连排
%   （0.5\linewidth−0.5\lxhang−0.25em）→单列 \lxopt；禁缩字号/负 kern/删标点凑宽。
%   本片实测降档：题1 两连排、题2/3 单列、题15 两连排、题16 单列（登记见值台账＋回执）。
% 解答书写区：E8~E14 七题均无小问，\liubai[16mm]{解答书写区}（默认 16mm；两问 24mm、
%   三问及以上 32mm、cap 40mm——本件全落默认档，逐题读数见值台账；纯题档吞 ansblock 不吞 \liubai）。
% 填空印答：E7 空位 \kongda{2}（详解版印答、纯题版退定宽空线，两档等形）。
% ============================================================
\documentclass[fontset=none]{ctexart}
\usepackage{qp-m3}
% —— 印面逐字符号路由补钉（承导学件母版；− 走 TNR、▱ NSC 子块；禁改 sty 保 md5） ——
\xeCJKDeclareCharClass{Default}{"2212}%
\setCJKfamilyfont{nscblk}[Path=C:/提示词/工作区/字替对照-0909/variantF/fonts/,BoldFont=NSC-w600.ttf]{NSC-w500.ttf}%
\xeCJKDeclareSubCJKBlock{nscblk}{"25B1}%
\newcommand{\pxparallelogram}{{\CJKfamily{nscblk}▱}}%
% —— 断行微调（承导学件母版实测档，三零门承重；\emergencystretch 不另设） ——
\xeCJKsetup{CJKglue={\hskip 0.10em plus 0.08em minus 0.05em}}%
% —— 练习件全线括线模（S3 波次方案 §四.3：导言一次置定、件内不切换） ——
\ansblockgrayfalse
% —— 页脚件名（qp-headfoot 复用入口） ——
\renewcommand{\qpzhangming}{第二章\quad 平面解析几何}%
\renewcommand{\qpjianming}{练习件}%

\begin{document}

\keshi{第1课时\quad 坐标法}
\begin{multicols}{2}
\emergencystretch=1em

% ============================================================
% 基础巩固（题1~7＝E1~E7：选择3＋填空4；题后紧跟答案制，全线括线 ansblock）
% ============================================================
\dangbadge{基}{础}{巩}{固}

\tihao{1}\zu{距离与中点公式直用×单选} \tieside{中档(知识点二)}已知四边形\(MNPQ\)的顶点\(M(1,1)\)，\(N(3,-1)\)，\(P(4,0)\)，\(Q(2,2)\)，则四边形\(MNPQ\)的形状为\nobreak(\kongwei)
\bindopt {\fontsize{10.5pt}{\qpTJgLeadLiti}\selectfont \makebox[\dimexpr0.5\linewidth-0.5\lxhang-0.25em\relax][l]{A．平行四边形}\makebox[\dimexpr0.5\linewidth-0.5\lxhang-0.25em\relax][l]{B．菱形}\par}
\bindopt {\fontsize{10.5pt}{\qpTJgLeadLiti}\selectfont \makebox[\dimexpr0.5\linewidth-0.5\lxhang-0.25em\relax][l]{C．梯形}\makebox[\dimexpr0.5\linewidth-0.5\lxhang-0.25em\relax][l]{D．矩形}\par}

\begin{ansblock}[2章-练-课时01-E1]
% ans:2章-练-课时01-E1
\ansitem{1}{D}
\ansnote{详解}{四边形\(MNPQ\)的对角线为\(MP\)与\(NQ\)．\(MP\)的中点为\(((1+4)/2,(1+0)/2)=(5/2,1/2)\)，\(NQ\)的中点为\(((3+2)/2,(-1+2)/2)=(5/2,1/2)\)，两对角线互相平分，故\(MNPQ\)是平行四边形．又 \(|MP|^{2}=(4-1)^{2}+(0-1)^{2}=10\)，\(|NQ|^{2}=(2-3)^{2}+(2+1)^{2}=10\)，对角线相等，故\pxparallelogram\(MNPQ\)为矩形．（若担心为正方形：\(|MN|^{2}=(3-1)^{2}+(-1-1)^{2}=8\)，\(|NP|^{2}=(4-3)^{2}+(0+1)^{2}=2\)，邻边不等，非正方形．）故选：D．}
\end{ansblock}

\tihao{2}\duoxuan\hspace{0.6em}\zu{距离与中点公式直用×多选} \tieside{中档(知识点二)}定义点\(P(x_0,y_0)\)到直线\(l\)：\(ax+by+c=0\)（\(a^{2}+b^{2}\neq 0\)）的有向距离为 \(d=(ax_0+by_0+c)/\sqrt{a^{2}+b^{2}}\)．已知点\(P_1\)，\(P_2\)到直线\(l\)的有向距离分别是\(d_1\)，\(d_2\)，则以下命题不正确的是\nobreak(\kongwei)

\lxopt{A．若\(d_1=d_2=1\)，则直线\(P_1P_2\)与直线\(l\)平行}

\lxopt{B．若\(d_1=1\)，\(d_2=-1\)，则直线\(P_1P_2\)与直线\(l\)垂直}

\lxopt{C．若\(d_1+d_2=0\)，则直线\(P_1P_2\)与直线\(l\)垂直}

\lxopt{D．若\(d_1\cdot d_2\leqslant 0\)，则直线\(P_1P_2\)与直线\(l\)相交}

\begin{ansblock}[2章-练-课时01-E2]
% ans:2章-练-课时01-E2
\ansitem{2}{BCD}
\ansnote{详解}{设\(P_1(x_1,y_1)\)，\(P_2(x_2,y_2)\)．A：\(d_1=d_2=1\)，则 \(ax_1+by_1+c=ax_2+by_2+c=\sqrt{a^{2}+b^{2}}\)，即 \(P_1\)、\(P_2\) 都在直线 \(ax+by+c-\sqrt{a^{2}+b^{2}}=0\) 上，该直线与\(l\)平行，故直线\(P_1P_2\parallel l\)，A正确；B：\(d_1=1\)，\(d_2=-1\) 只说明\(P_1\)、\(P_2\)在\(l\)的两侧且到\(l\)的距离相等，线段\(P_1P_2\)的中点落在\(l\)上，但\(P_1P_2\)与\(l\)不一定垂直，B错误；C：\(d_1+d_2=0\) 包含\(d_1=d_2=0\)的情形，此时\(P_1\)、\(P_2\)都在\(l\)上，直线\(P_1P_2\)与\(l\)重合，谈不上垂直，C错误；D：\(d_1\cdot d_2\leqslant 0\) 表示\(P_1\)、\(P_2\)分居\(l\)两侧或至少一点在\(l\)上，此时直线\(P_1P_2\)与\(l\)相交或重合，D把“重合”漏掉，D错误．故选：BCD．}
\end{ansblock}

\tihao{3}\duoxuan\hspace{0.6em}\zu{距离与中点公式直用×多选} \tieside{中档(知识点二)}在平面直角坐标系中，已知点\(A(x_1,y_1)\)、\(B(x_2,y_2)\)，定义 \(d(A,B)=|x_1-x_2|+|y_1-y_2|\) 为两点\(A\)，\(B\)的“折线距离”；又设点\(P\)及直线\(l\)上任意一点\(Q\)，称\(d(P,Q)\)的最小值为点\(P\)到直线\(l\)的“折线距离”，记作\(d(P,l)\)．则下列说法正确的是\nobreak(\kongwei)

\lxopt{A．对任意的两点\(A\)，\(B\)，都有 \(d(A,B)\geqslant|AB|\)}

\lxopt{B．对任意三点\(A\)、\(B\)、\(C\)，都有 \(d(A,C)+d(B,C)\geqslant d(A,B)\)}

\lxopt{C．已知点\(P(3,1)\)和直线\(l\)：\(2x-y-1=0\)，则 \(d(P,l)=4\sqrt{5}/5\)}

\lxopt{D．已知点\(O(0,0)\)，动点\(P(x,y)\)满足\(d(P,O)=1\)，则动点\(P\)的轨迹围成平面图形的面积是2}

\begin{ansblock}[2章-练-课时01-E3]
% ans:2章-练-课时01-E3
\ansitem{3}{ABD}
\ansnote{详解}{A：\((|x_1-x_2|+|y_1-y_2|)^{2}=(x_1-x_2)^{2}+2|x_1-x_2||y_1-y_2|+(y_1-y_2)^{2}\geqslant(x_1-x_2)^{2}+(y_1-y_2)^{2}\)，开方即得 \(d(A,B)\geqslant|AB|\)，A正确；B：由绝对值三角不等式，\(|x_1-x_2|+|x_2-x_3|\geqslant|x_1-x_3|\)，\(|y_1-y_2|+|y_2-y_3|\geqslant|y_1-y_3|\)，两式相加即得 \(d(A,B)+d(B,C)\geqslant d(A,C)\)（将\(A\)、\(B\)、\(C\)重排即题设形式），B正确；C：设\(Q(x,2x-1)\)是\(l\)上任一点，\(d(P,Q)=|x-3|+|2x-2|=|x-3|+|x-1|+|x-1|\geqslant|(x-3)-(x-1)|+0=2\)，当且仅当\(x=1\)时等号成立，故 \(d(P,l)=2\neq 4\sqrt{5}/5\)，C错误；D：\(d(P,O)=|x|+|y|=1\) 的图象是以\((1,0)\)、\((0,1)\)、\((-1,0)\)、\((0,-1)\)为顶点的正方形，对角线长均为2，面积\(=(1/2)\times 2\times 2=2\)，D正确．故选：ABD．}
\end{ansblock}

\tihao{4}\zu{距离与中点公式直用×填空} \tieside{简单(知识点二)}已知\(A(3,0)\)，\(B(0,-4)\)，求这两点之间的距离以及它们的中点坐标．

\begin{ansblock}[2章-练-课时01-E4]
% ans:2章-练-课时01-E4
\ansitem{4}{|AB|＝5；中点坐标为(3/2,−2)．}
\ansnote{详解}{\(|AB|=\sqrt{(0-3)^{2}+(-4-0)^{2}}=\sqrt{9+16}=5\)；中点坐标为\(((3+0)/2,(0-4)/2)=(3/2,-2)\)．}
\end{ansblock}

\tihao{5}\zu{距离与中点公式直用×填空} \tieside{简单(知识点一)}已知数轴上\(A(-1)\)，\(B(2)\)，且\(A\)关于\(B\)的对称点为\(C\)，求\(C\)的坐标．

\begin{ansblock}[2章-练-课时01-E5]
% ans:2章-练-课时01-E5
\ansitem{5}{C的坐标为5．}
\ansnote{详解}{\(B\)是线段\(AC\)的中点，设\(C(x)\)，则 \((-1+x)/2=2\)，解得\(x=5\)，即\(C\)的坐标为5．}
\end{ansblock}

\tihao{6}\zu{距离与中点公式直用×填空} \tieside{简单(知识点一)}已知数轴上的点\(P\)到\(A(-9)\)的距离是它到\(B(-3)\)的距离的2倍，求点\(P\)的坐标．

\begin{ansblock}[2章-练-课时01-E6]
% ans:2章-练-课时01-E6
\ansitem{6}{P的坐标为3或−5．}
\ansnote{详解}{设\(P(x)\)，依题意 \(|x+9|=2|x+3|\)．两边均非负，平方得 \((x+9)^{2}=4(x+3)^{2}\)，即 \(x^{2}+18x+81=4x^{2}+24x+36\)，整理得 \(3x^{2}+6x-45=0\)，即 \(x^{2}+2x-15=0\)，解得\(x=3\)或\(x=-5\)．检验：\(x=3\)时\(|3+9|=12=2|3+3|\)；\(x=-5\)时\(|-5+9|=4=2|-5+3|\)．所以点\(P\)的坐标为3或\(-5\)．}
\end{ansblock}

\tihao{7}\zu{距离与中点公式直用×填空} \tieside{中档(知识点二)}定义点到曲线的距离为该点与曲线上所有点之间距离的最小值，则点\(P(0,1)\)到曲线 \(x-xy+2=0\) 的距离为\kongda{2}．

\begin{ansblock}[2章-练-课时01-E7]
% ans:2章-练-课时01-E7
\ansitem{7}{2}
\ansnote{详解}{曲线方程中\(x=0\)时\(2=0\)不成立，故\(x\neq 0\)，且 \(y=1+2/x\)．设曲线上任意一点\(A(x,1+2/x)\)，则\(|AP|=\sqrt{x^{2}+(1+2/x-1)^{2}}=\sqrt{x^{2}+4/x^{2}}\)．由基本不等式 \(x^{2}+4/x^{2}\geqslant 2\sqrt{x^{2}\cdot 4/x^{2}}=4\)，当且仅当 \(x^{2}=4/x^{2}\)，即\(x=\pm\sqrt{2}\) 时等号成立，此时\(|AP|\)最小值为\(\sqrt{4}=2\)．故所求距离为2．}
\end{ansblock}

% ============================================================
% 综合提升（题8~14＝E8~E14：解答7，书写区 \liubai[16mm]（无小问默认档）＋题后括线 ansblock）
% ============================================================
\dangbadge{综}{合}{提}{升}

\tihao{8}\zu{距离与中点公式直用×解答} \tieside{中档(知识点二)}已知\(A(-1,3)\)，\(B(3,3)\)，\(C(1,2\sqrt{3}+3)\)，证明\(\triangle ABC\)是等边三角形．
\liubai[16mm]{解答书写区}

\begin{ansblock}[2章-练-课时01-E8]
% ans:2章-练-课时01-E8
\ansitem{8}{证明见详解．}
\ansnote{详解}{由两点间距离公式：\par\noindent \(|AB|=\sqrt{(3+1)^{2}+(3-3)^{2}}=\sqrt{16}=4\)；\par\noindent \(|AC|=\sqrt{(1+1)^{2}+(2\sqrt{3}+3-3)^{2}}=\sqrt{4+12}=4\)；\par\noindent \(|BC|=\sqrt{(1-3)^{2}+(2\sqrt{3}+3-3)^{2}}=\sqrt{4+12}=4\)．\par\noindent 所以\(|AB|=|AC|=|BC|\)，故\(\triangle ABC\)是等边三角形．}
\end{ansblock}

\tihao{9}\zu{距离与中点公式直用×解答} \tieside{中档(知识点二)}已知\(\triangle ABC\)的三个顶点分别为\(A(2,1)\)，\(B(-2,3)\)，\(C(0,-1)\)，求\(\triangle ABC\)三条中线的长．
\liubai[16mm]{解答书写区}

\begin{ansblock}[2章-练-课时01-E9]
% ans:2章-练-课时01-E9
\ansitem{9}{AB边上的中线长3；BC边上的中线长3；CA边上的中线长3√2．}
\ansnote{详解}{设\(AB\)、\(BC\)、\(CA\)的中点分别为\(M\)、\(N\)、\(P\)．\(M((2-2)/2,(1+3)/2)=M(0,2)\)，中线\(|CM|=\sqrt{(0-0)^{2}+(2+1)^{2}}=3\)；\(N((-2+0)/2,(3-1)/2)=N(-1,1)\)，中线\(|AN|=\sqrt{(-1-2)^{2}+(1-1)^{2}}=3\)；\(P((0+2)/2,(-1+1)/2)=P(1,0)\)，中线\(|BP|=\sqrt{(1+2)^{2}+(0-3)^{2}}=\sqrt{18}=3\sqrt{2}\)．所以三条中线的长分别为3、3、\(3\sqrt{2}\)．}
\end{ansblock}

\tihao{10}\zu{距离与中点公式直用×解答} \tieside{中档(知识点三)}已知\(A(3,1)\)，\(B(-2,2)\)，在\(y\)轴上的点\(P\)满足\(PA\perp PB\)，求\(P\)的坐标．
\liubai[16mm]{解答书写区}

\begin{ansblock}[2章-练-课时01-E10]
% ans:2章-练-课时01-E10
\ansitem{10}{P(0,4)或P(0,−1)．}
\ansnote{详解}{设\(P(0,y)\)．\(PA\perp PB\) 即\(\triangle PAB\)是以\(AB\)为斜边的直角三角形，\(|PA|^{2}+|PB|^{2}=|AB|^{2}\)．\(|PA|^{2}=9+(y-1)^{2}\)，\(|PB|^{2}=4+(y-2)^{2}\)，\(|AB|^{2}=(-2-3)^{2}+(2-1)^{2}=26\)，代入：\(9+y^{2}-2y+1+4+y^{2}-4y+4=26\)，整理得 \(2y^{2}-6y-8=0\)，即 \(y^{2}-3y-4=0\)，解得\(y=4\)或\(y=-1\)．所以点\(P\)的坐标为\((0,4)\)或\((0,-1)\)．}
\end{ansblock}

\tihao{11}\zu{坐标法证明与存在性讨论×解答} \tieside{中档(知识点三)}已知\(ABCD\)是一个长方形，且\(M\)是\(ABCD\)所在平面上任意一个点，求证：\(|AM|^{2}+|CM|^{2}=|BM|^{2}+|DM|^{2}\)．
\liubai[16mm]{解答书写区}

\begin{ansblock}[2章-练-课时01-E11]
% ans:2章-练-课时01-E11
\ansitem{11}{证明见详解．}
\ansnote{详解}{以长方形的顶点\(A\)为原点、\(AB\)所在直线为\(x\)轴建立平面直角坐标系．设\(AB=a\)，\(AD=b\)，则\(A(0,0)\)，\(B(a,0)\)，\(C(a,b)\)，\(D(0,b)\)．再设\(M(x,y)\)．\(|AM|^{2}+|CM|^{2}=x^{2}+y^{2}+(x-a)^{2}+(y-b)^{2}\)；\(|BM|^{2}+|DM|^{2}=(x-a)^{2}+y^{2}+x^{2}+(y-b)^{2}\)．两式展开后均为 \(x^{2}+y^{2}+(x-a)^{2}+(y-b)^{2}\)，恒相等．所以 \(|AM|^{2}+|CM|^{2}=|BM|^{2}+|DM|^{2}\)，即对平面上任意一点\(M\)结论都成立．}
\end{ansblock}

\tihao{12}\zu{坐标法证明与存在性讨论×解答} \tieside{中档(知识点三)}已知函数\(y=f(x)\)的图象与函数\(y=x^{2}+1\)的图象关于点\(M(2,0)\)对称，求\(f(x)\)的解析式．
\liubai[16mm]{解答书写区}

\begin{ansblock}[2章-练-课时01-E12]
% ans:2章-练-课时01-E12
\ansitem{12}{f(x)＝−(x−4)²−1．}
\ansnote{详解}{在\(y=f(x)\)的图象上任取一点\(P(x,y)\)，它关于点\(M(2,0)\)的对称点为 \(P'(4-x,-y)\)．由对称性，\(P'\)在函数\(y=x^{2}+1\)的图象上，所以 \(-y=(4-x)^{2}+1\)，即 \(y=-(x-4)^{2}-1\)．故 \(f(x)=-(x-4)^{2}-1\)．}
\end{ansblock}

\tihao{13}\zu{坐标法证明与存在性讨论×解答} \tieside{中档(知识点三)}已知正方形\(ABCD\)的边长为4，若\(E\)是\(BC\)的中点，\(F\)是\(CD\)的中点，求证：\(BF\perp AE\)．
\liubai[16mm]{解答书写区}

\begin{ansblock}[2章-练-课时01-E13]
% ans:2章-练-课时01-E13
\ansitem{13}{证明见详解．}
\ansnote{详解}{以\(A\)为原点、\(AB\)所在直线为\(x\)轴建立平面直角坐标系，则\(A(0,0)\)，\(B(4,0)\)，\(C(4,4)\)，\(D(0,4)\)；由\(E\)、\(F\)分别是\(BC\)、\(CD\)的中点得 \(E(4,2)\)，\(F(2,4)\)．于是 \(\overrightarrow{AE}=(4,2)\)，\(\overrightarrow{BF}=(2-4,4-0)=(-2,4)\)．\(\overrightarrow{AE}\cdot\overrightarrow{BF}=4\times(-2)+2\times 4=-8+8=0\)，所以 \(\overrightarrow{AE}\perp\overrightarrow{BF}\)，即 \(BF\perp AE\)．}
\end{ansblock}

\tihao{14}\zu{坐标法证明与存在性讨论×解答} \tieside{中档(知识点三)}求证：平行四边形两条对角线的平方和等于四条边的平方和．
\liubai[16mm]{解答书写区}

\begin{ansblock}[2章-练-课时01-E14]
% ans:2章-练-课时01-E14
\ansitem{14}{证明见详解．}
\ansnote{详解}{以平行四边形相邻两边为坐标轴方向建立平面直角坐标系：取顶点\(A\)为原点、\(AB\)所在直线为\(x\)轴．设\(B(a,0)\)，\(D(b,c)\)，则由平行四边形性质知 \(C(a+b,c)\)．两条对角线：\(|AC|^{2}=(a+b)^{2}+c^{2}\)，\(|BD|^{2}=(b-a)^{2}+c^{2}\)，所以 \(|AC|^{2}+|BD|^{2}=\left(a^{2}+2ab+b^{2}+c^{2}\right)+\left(a^{2}-2ab+b^{2}+c^{2}\right)=2a^{2}+2b^{2}+2c^{2}\)．四条边：\(|AB|^{2}=a^{2}\)，\(|CD|^{2}=a^{2}\)，\(|AD|^{2}=b^{2}+c^{2}\)，\(|BC|^{2}=b^{2}+c^{2}\)，所以 \(|AB|^{2}+|BC|^{2}+|CD|^{2}+|DA|^{2}=2a^{2}+2b^{2}+2c^{2}\)．两式相同，故平行四边形两条对角线的平方和等于四条边的平方和．}
\end{ansblock}

% ============================================================
% 思维探索（题15~16＝E15~E16：单选2，0913命制入槽新命对，居末档照库序）
% ============================================================
\dangbadge{思}{维}{探}{索}

\tihao{15}\zu{坐标法证明与存在性讨论×单选} \tieside{简单(知识点三)}在 Rt\(\triangle ABC\) 中，\(\angle C=90^\circ\)，以 \(C\) 为原点、\(CB\)、\(CA\) 所在直线分别为 \(x\) 轴、\(y\) 轴建立平面直角坐标系，设 \(B(b,0)\)、\(A(0,a)\)（\(a>0\)，\(b>0\)），\(M\) 为斜边 \(AB\) 的中点，则下列结论正确的是\nobreak(\kongwei)
\bindopt {\fontsize{10.5pt}{\qpTJgLeadLiti}\selectfont \makebox[\dimexpr0.5\linewidth-0.5\lxhang-0.25em\relax][l]{A．\(|CM|=|AB|\)}\makebox[\dimexpr0.5\linewidth-0.5\lxhang-0.25em\relax][l]{B．\(|CM|=|AB|/2\)}\par}
\bindopt {\fontsize{10.5pt}{\qpTJgLeadLiti}\selectfont \makebox[\dimexpr0.5\linewidth-0.5\lxhang-0.25em\relax][l]{C．\(|CM|=(a+b)/2\)}\makebox[\dimexpr0.5\linewidth-0.5\lxhang-0.25em\relax][l]{D．\(|CM|=(a^{2}+b^{2})/2\)}\par}

\begin{ansblock}[2章-练-课时01-E15]
% ans:2章-练-课时01-E15
\ansitem{15}{B}
\ansnote{详解}{由中点坐标公式，\(M((b+0)/2,(0+a)/2)\)，即 \(M(b/2,a/2)\)。由两点间距离公式：\(|CM|=\sqrt{(b/2)^{2}+(a/2)^{2}}=\sqrt{a^{2}+b^{2}}/2\)；\(|AB|=\sqrt{(b-0)^{2}+(0-a)^{2}}=\sqrt{a^{2}+b^{2}}\)。故 \(|CM|=|AB|/2\)，选 B。这就用坐标法证明了：直角三角形斜边上的中线等于斜边的一半。}
\end{ansblock}

\tihao{16}\zu{坐标法证明与存在性讨论×单选} \tieside{简单(知识点三)}在\(\triangle ABC\) 中，以 \(B\) 为原点、\(BC\) 所在直线为 \(x\) 轴建立平面直角坐标系，设 \(C(a,0)\)、\(A(b,c)\)（\(a>0\)，\(c>0\)）。\(M\)、\(N\) 分别为 \(AB\)、\(AC\) 的中点，则下列关系成立的是\nobreak(\kongwei)

\lxopt{A．\(MN\parallel BC\)，且 \(|MN|=|BC|\)}

\lxopt{B．\(MN\parallel BC\)，且 \(|MN|=|BC|/2\)}

\lxopt{C．\(MN\perp BC\)，且 \(|MN|=|BC|/2\)}

\lxopt{D．\(MN\) 与 \(BC\) 不平行，且 \(|MN|=|BC|/2\)}

\begin{ansblock}[2章-练-课时01-E16]
% ans:2章-练-课时01-E16
\ansitem{16}{B}
\ansnote{详解}{由中点坐标公式，\(M(b/2,c/2)\)，\(N((a+b)/2,c/2)\)。两点纵坐标相等且 \(|MN|=a/2>0\)（\(M\neq N\)），故 \(MN\parallel x\) 轴；而 \(BC\) 在 \(x\) 轴上，所以 \(MN\parallel BC\)。又 \(|MN|=|(a+b)/2-b/2|=a/2\)，\(|BC|=a\)，故 \(|MN|=|BC|/2\)。选 B。这就用坐标法证明了：三角形中位线平行于第三边，且等于第三边的一半。}
\end{ansblock}

% ---- 尾块（\tailfill 置 \end{multicols} 前；无参＝内容尾随框，每件恰 1 处） ----
\tailfill

\end{multicols}
\end{document}
